% GENERATED FILE. Edit canonical lessons, then run npm run generate:books.
% canonical-source-hash: fnv1a64:1f810048d0f8e523
% canonical-lessons: KA-C385-prabhava, KA-C385-badhisu, KA-C385-untagu, KA-C385-pravrtti, KA-C385-mula

\chapter{Causes and Effects}
\label{ch:ka-causes-and-effects}
\hlchaptermodality{385}

\begin{chapteropening}
\textbf{By the end of this chapter:} \emph{I can say influence, to affect badly, to come about, a tendency and a source.}

Complete the last lesson of chapter 385: I can say influence, to affect badly, to come about, a tendency and a source.
\end{chapteropening}

\section[\texorpdfstring{\kn{ಪ್ರಭಾವ} (prabhāva)}{ಪ್ರಭಾವ (prabhāva)}]{\kn{ಪ್ರಭಾವ} (prabhāva) — influence}

\label{lesson:KA-C385-prabhava}



\begin{quote}
Before the new one: say the Kannada for an aspect, then the Kannada for a feature.
\end{quote}

\subsection*{You'll want to know: \kn{ಪ್ರಭಾವ}}
\textbf{\kn{ಪ್ರಭಾವ}} — \emph{prabhāva} — \textquotedblleft{}influence\textquotedblright{}.

\begin{quote}
\kn{ಅವನ} \kn{ಮೇಲೆ} \kn{ಅಜ್ಜನ} \kn{ಪ್ರಭಾವ} \kn{ಹೆಚ್ಚು}.
\end{quote}

\begin{quote}
\emph{avana mēle ajjana prabhāva heccu.}
\end{quote}

\begin{quote}
His grandfather's influence on him is strong.
\end{quote}

\kn{ಮೇಲೆ} marks who is influenced.

\begin{cousinweb}
\textbf{\kn{ಪ್ರಭಾವ}} (\emph{prabhāva}) is borrowed from Sanskrit \emph{prabhāva}, \textquotedblleft{}might, power, majesty\textquotedblright{} (Monier-Williams).
\end{cousinweb}

\begin{grammarlens}[title={What you've built}]
One more word for giving reasons and opinions.
\end{grammarlens}

\subsection*{Your turn}
\begin{itemize}
\raggedright
  \item \emph{Say it:} \emph{prabhāva}
  \item \emph{Say it:} \emph{prabhāva}, once more
  \item \emph{Say it:} \emph{avana mēle ajjana prabhāva heccu.} — His grandfather's influence on him is strong.
  \item \emph{Recall:} say the Kannada for an aspect, then the Kannada for a feature, then say \emph{prabhāva} again
\end{itemize}

\begin{tcolorbox}[breakable,colback=teal!4,colframe=teal!35!black,title={Before you move on}]
What does \kn{ಪ್ರಭಾವ} mean? (\textquotedblleft{}Influence\textquotedblright{}.) Say it once more.
\end{tcolorbox}
\section[\texorpdfstring{\kn{ಬಾಧಿಸು} (bādhisu)}{ಬಾಧಿಸು (bādhisu)}]{\kn{ಬಾಧಿಸು} (bādhisu) — to affect badly}

\label{lesson:KA-C385-badhisu}



\begin{quote}
Before the new one: say the Kannada for a feature, then the Kannada for influence.
\end{quote}

\subsection*{You'll want to know: \kn{ಬಾಧಿಸು}}
\textbf{\kn{ಬಾಧಿಸು}} — \emph{bādhisu} — \textquotedblleft{}to affect badly\textquotedblright{}.

\begin{quote}
\kn{ಬೆಲೆ} \kn{ಏರಿಕೆ} \kn{ಬಡವರನ್ನು} \kn{ಹೆಚ್ಚು} \kn{ಬಾಧಿಸಿತು}.
\end{quote}

\begin{quote}
\emph{bele ērike baḍavarannu heccu bādhisitu.}
\end{quote}

\begin{quote}
The rise in prices hit the poor hardest.
\end{quote}

Past (it …): \kn{ಬಾಧಿಸಿತು}, \emph{bādhisitu}.

Who is affected takes -\kn{ಅನ್ನು}. \kn{ಬಾಧೆ}, distress.

\begin{cousinweb}
\textbf{\kn{ಬಾಧಿಸು}} (\emph{bādhisu}) is built on Sanskrit \emph{bādhana}, \textquotedblleft{}oppressing, harassing\textquotedblright{} (Monier-Williams), with the Kannada verb ending -\emph{isu}.
\end{cousinweb}

\begin{grammarlens}[title={What you've built}]
One more word for giving reasons and opinions.
\end{grammarlens}

\subsection*{Your turn}
\begin{itemize}
\raggedright
  \item \emph{Say it:} \emph{bādhisu}
  \item \emph{Say it:} \emph{bādhisu}, once more
  \item \emph{Say it:} \emph{bele ērike baḍavarannu heccu bādhisitu.} — The rise in prices hit the poor hardest.
  \item \emph{Recall:} say the Kannada for a feature, then the Kannada for influence, then say \emph{bādhisu} again
\end{itemize}

\begin{tcolorbox}[breakable,colback=teal!4,colframe=teal!35!black,title={Before you move on}]
What does \kn{ಬಾಧಿಸು} mean? (\textquotedblleft{}To affect badly\textquotedblright{}.) Say it once more.
\end{tcolorbox}
\section[\texorpdfstring{\kn{ಉಂಟಾಗು} (uṇṭāgu)}{ಉಂಟಾಗು (uṇṭāgu)}]{\kn{ಉಂಟಾಗು} (uṇṭāgu) — to come about; to arise}

\label{lesson:KA-C385-untagu}



\begin{quote}
Before the new one: say the Kannada for influence, then the Kannada for to affect badly.
\end{quote}

\subsection*{You'll want to know: \kn{ಉಂಟಾಗು}}
\textbf{\kn{ಉಂಟಾಗು}} — \emph{uṇṭāgu} — \textquotedblleft{}to come about; to arise\textquotedblright{}.

\begin{quote}
\kn{ಮಳೆಯಿಂದ} \kn{ರಸ್ತೆಯಲ್ಲಿ} \kn{ತೊಂದರೆ} \kn{ಉಂಟಾಯಿತು}.
\end{quote}

\begin{quote}
\emph{maḷeyinda rasteyalli tondare uṇṭāyitu.}
\end{quote}

\begin{quote}
The rain caused trouble on the road.
\end{quote}

Past (it …): \kn{ಉಂಟಾಯಿತು}, \emph{uṇṭāyitu}.

\kn{ಉಂಟು}, there is, + \kn{ಆಗು}. \kn{ಉಂಟುಮಾಡು}, to cause.

\begin{cousinweb}
\textbf{\kn{ಉಂಟಾಗು}} (\emph{uṇṭāgu}) is a native Dravidian compound: \emph{uṇṭu}, \textquotedblleft{}there is\textquotedblright{} (DEDR 697), and \emph{āgu}, \textquotedblleft{}to become\textquotedblright{} (DEDR 333).
\end{cousinweb}

\begin{grammarlens}[title={What you've built}]
One more word for giving reasons and opinions.
\end{grammarlens}

\subsection*{Your turn}
\begin{itemize}
\raggedright
  \item \emph{Say it:} \emph{uṇṭāgu}
  \item \emph{Say it:} \emph{uṇṭāgu}, once more
  \item \emph{Say it:} \emph{maḷeyinda rasteyalli tondare uṇṭāyitu.} — The rain caused trouble on the road.
  \item \emph{Recall:} say the Kannada for influence, then the Kannada for to affect badly, then say \emph{uṇṭāgu} again
\end{itemize}

\begin{tcolorbox}[breakable,colback=teal!4,colframe=teal!35!black,title={Before you move on}]
What does \kn{ಉಂಟಾಗು} mean? (\textquotedblleft{}To come about; to arise\textquotedblright{}.) Say it once more.
\end{tcolorbox}
\section[\texorpdfstring{\kn{ಪ್ರವೃತ್ತಿ} (pravṛtti)}{ಪ್ರವೃತ್ತಿ (pravṛtti)}]{\kn{ಪ್ರವೃತ್ತಿ} (pravṛtti) — a tendency; a trend}

\label{lesson:KA-C385-pravrtti}



\begin{quote}
Before the new one: say the Kannada for to affect badly, then the Kannada for to come about.
\end{quote}

\subsection*{You'll want to know: \kn{ಪ್ರವೃತ್ತಿ}}
\textbf{\kn{ಪ್ರವೃತ್ತಿ}} — \emph{pravṛtti} — \textquotedblleft{}a tendency; a trend\textquotedblright{}.

\begin{quote}
\kn{ಯುವಕರಲ್ಲಿ} \kn{ನಗರಕ್ಕೆ} \kn{ಹೋಗುವ} \kn{ಪ್ರವೃತ್ತಿ} \kn{ಹೆಚ್ಚಾಗಿದೆ}.
\end{quote}

\begin{quote}
\emph{yuvakaralli nagarakke hōguva pravṛtti heccāgide.}
\end{quote}

\begin{quote}
Among young people the trend of moving to the city has grown.
\end{quote}

A verb in -\kn{ುವ} comes before it: \kn{ಹೋಗುವ} \kn{ಪ್ರವೃತ್ತಿ}.

\begin{cousinweb}
\textbf{\kn{ಪ್ರವೃತ್ತಿ}} (\emph{pravṛtti}) is borrowed from Sanskrit \emph{pravṛtti}, \textquotedblleft{}coming forth, activity, a course of action\textquotedblright{} (Monier-Williams).
\end{cousinweb}

\begin{grammarlens}[title={What you've built}]
One more word for giving reasons and opinions.
\end{grammarlens}

\subsection*{Your turn}
\begin{itemize}
\raggedright
  \item \emph{Say it:} \emph{pravṛtti}
  \item \emph{Say it:} \emph{pravṛtti}, once more
  \item \emph{Say it:} \emph{yuvakaralli nagarakke hōguva pravṛtti heccāgide.} — Among young people the trend of moving to the city has grown.
  \item \emph{Recall:} say the Kannada for to affect badly, then the Kannada for to come about, then say \emph{pravṛtti} again
\end{itemize}

\begin{tcolorbox}[breakable,colback=teal!4,colframe=teal!35!black,title={Before you move on}]
What does \kn{ಪ್ರವೃತ್ತಿ} mean? (\textquotedblleft{}A tendency; a trend\textquotedblright{}.) Say it once more.
\end{tcolorbox}
\section[\texorpdfstring{\kn{ಮೂಲ} (mūla)}{ಮೂಲ (mūla)}]{\kn{ಮೂಲ} (mūla) — a source; a root}

\label{lesson:KA-C385-mula}



\begin{quote}
Before the new one: say the Kannada for to come about, then the Kannada for a tendency.
\end{quote}

\subsection*{You'll want to know: \kn{ಮೂಲ}}
\textbf{\kn{ಮೂಲ}} — \emph{mūla} — \textquotedblleft{}a source; a root\textquotedblright{}.

\begin{quote}
\kn{ಈ} \kn{ನದಿಯ} \kn{ಮೂಲ} \kn{ಬೆಟ್ಟದಲ್ಲಿದೆ}.
\end{quote}

\begin{quote}
\emph{ī nadiya mūla beṭṭadallide.}
\end{quote}

\begin{quote}
This river's source is in the hills.
\end{quote}

\kn{ಮೂಲತಃ}, originally.

\begin{cousinweb}
\textbf{\kn{ಮೂಲ}} (\emph{mūla}) is borrowed from Sanskrit \emph{mūla}, \textquotedblleft{}a root, a base, an origin\textquotedblright{} (CDIAL 10250).
\end{cousinweb}

\begin{grammarlens}[title={What you've built}]
One more word for giving reasons and opinions.
\end{grammarlens}

\subsection*{Your turn}
\begin{itemize}
\raggedright
  \item \emph{Say it:} \emph{mūla}
  \item \emph{Say it:} \emph{mūla}, once more
  \item \emph{Say it:} \emph{ī nadiya mūla beṭṭadallide.} — This river's source is in the hills.
  \item \emph{Recall:} say the Kannada for to come about, then the Kannada for a tendency, then say \emph{mūla} again
\end{itemize}

\begin{tcolorbox}[breakable,colback=teal!4,colframe=teal!35!black,title={Before you move on}]
What does \kn{ಮೂಲ} mean? (\textquotedblleft{}A source; a root\textquotedblright{}.) Say it once more.
\end{tcolorbox}
